% Brief (tikz-diagram-craft). Figure 4.1, The Legible Swarm.
% Reader question: what can actually happen when three independently launched
%   agents edit the same developer settings file from the same revision?
% Claim: an ungoverned overlap can converge, break loudly, or remain valid and
%   fail later; visible claims, locks, and downstream evidence are what let an
%   operator distinguish those cases.
% Evidence: one concrete .vscode/settings.json base revision; three proposed
%   key edits; a pointer claim, a whole-file lock, and a second pointer claim;
%   three possible resolutions and their first observable consequences
%   [internal, pedagogical counterexample].
% Must distinguish: launch source from edit intent; pointer claim from file
%   lock; syntactic failure from semantic failure; test-time observation from
%   downstream consequence.
% Grammar chosen: three aligned worked traces from one shared input. Equal
%   rows make the operation and first observable output directly comparable;
%   the delayed row alone continues to its later release failure.
% Grammar rejected: actor bubbles around settings.json. A membership graph
%   would show who touched the file but erase read revision, order, overlap,
%   and the delay between a green mutation score and the downstream failure.
% Five-second test: three agents read a1c9; one overlap is safe, one breaks now,
%   and one stays green until a later release consumes the hidden fixture.
\begin{figure}[H]
\centering
\pdfigurehue{pdteal}%
\begin{tikzpicture}[pd figure,x=1cm,y=1cm]
  \tikzset{
    actor/.style={pd artifact,text width=2.20cm,minimum height=1.40cm},
    operation/.style={pd state,text width=3.15cm,minimum height=1.35cm},
    result/.style={text width=3.50cm,minimum height=1.35cm}
  }

  % Every trace starts from this exact revision.
  \node[pd artifact,text width=10.35cm,align=left] (base) at (5.45,7.20)
    {\textbf{shared input}\quad \texttt{.vscode/settings.json @ a1c9}\\
     \texttt{formatOnSave=false}\quad \texttt{tsdk=node\_modules/...}\\
     \texttt{search.exclude=\{"**/dist": true\}}};
  \node[pd label,anchor=north] at (5.45,6.50)
    {all three agents read \texttt{a1c9} before any write lands};

  \node[actor] (A) at (1.30,4.82)
    {\textbf{Agent A}\\02:00 format cron\\\pdfigsub pointer claim};
  \node[operation] (AO) at (5.20,4.82)
    {set \texttt{formatOnSave}\\\texttt{false -> true}};
  \node[pd ready state,result] (AR) at (9.30,4.82)
    {\textbf{benign}\\keywise merge\\valid JSON; all keys};

  \node[actor] (B) at (1.30,2.62)
    {\textbf{Agent B}\\02:01 PnP prototype\\\pdfigsub whole-file lock};
  \node[operation] (BO) at (5.20,2.62)
    {replace \texttt{tsdk}\\by raw text splice};
  \node[pd breach state,result] (BR) at (9.30,2.62)
    {\textbf{immediate}\\comma lost\\parse fails at write};

  \node[actor] (C) at (1.30,0.42)
    {\textbf{Agent C}\\02:02 search cron\\\pdfigsub pointer claim};
  \node[operation] (CO) at (5.20,0.42)
    {add \texttt{**/fixtures}\\to \texttt{search.exclude}};
  \node[pd warn state,result] (CR) at (9.30,0.42)
    {\textbf{delayed}\\valid JSON; tests 98\%\\settings path untested};
  \node[pd breach state,result] (CL) at (9.30,-2.10)
    {\textbf{next release}\\\texttt{fixtures/v1} unseen\\replay fails};

  \begin{scope}[on background layer]
    \draw[pd focus arrow] (A.east) -- (AO.west);
    \draw[pd ready arrow] (AO.east) -- (AR.west);
    \draw[pd focus arrow] (B.east) -- (BO.west);
    \draw[pd breach arrow] (BO.east) -- (BR.west);
    \draw[pd focus arrow] (C.east) -- (CO.west);
    \draw[pd warn arrow] (CO.east) -- (CR.west);
    \draw[pd breach arrow] (CR.south) -- (CL.north);
  \end{scope}
\end{tikzpicture}
\caption{Three worked traces start from one \texttt{settings.json} revision.
A keywise merge is benign; a raw splice fails immediately; a valid exclusion
stays green until release replay exposes the hidden fixture. [internal,
pedagogical counterexample]}
\label{fig:state-of-nature}
\end{figure}
